\section{A rational-oracle lower bound for the framework}
\label{simp:lower-bound}

The benchmark in this section is unrestricted Nash social welfare.
The algorithm must first select $n$ reserved items $H$ by a
maximum-product singleton matching and must finally give exactly
one item of $H$ to every agent.  It may otherwise coordinate the
residual allocation and final matching arbitrarily.

\begin{theorem}\label{thm:simp-framework-lower-bound}
Let $R_\star=4^{1/4}(1-e^{-1})^{-3/4}$.  For every fixed
$0<\varepsilon<R_\star-1$, there are constants $n_0,c>0$ and,
on arbitrarily large item sets of size $m$, distributions of
normalized, nonnegative, monotone submodular valuations such that:
the number of unweighted agents is $n_0$; all answers are rational
with $O_\varepsilon(\log m)$ bits; every maximum-product singleton
matching selects the same specified set $H$; and any adaptive
randomized algorithm using at most $\exp(cm)$ queries and always
returning one reserved item per agent fails to attain factor
$R_\star-\varepsilon$ with probability tending to one.
\end{theorem}

\subsection{A smooth pair with exact agreement near the diagonal}

Fix $k\ge2$, let $\rho_k=1-(1-1/k)^k$, and, for
$x\in[0,1]^k$, write
\[
 \bar x=\frac1k\sum_jx_j,\quad
 F(x)=1-\prod_j(1-x_j),\quad G(x)=1-(1-\bar x)^k.
\]
Both functions are increasing and have nonpositive second partials,
including diagonal partials.  Arithmetic--geometric mean gives $F\ge G$.
We construct increasing continuous diminishing-returns functions
$f,g$ with, for any prescribed $\zeta>0$,
\begin{equation}\label{simp:masked-properties}
 \begin{gathered}
 f(0)=g(0)=0,\qquad 0\le\partial_jf,\partial_jg\le1+C,
 \qquad C<\zeta,\qquad f(e_j)\ge1-\zeta,\\
 \max_{p,q}|x_p-x_q|\le\eta
       \ \Longrightarrow\ f(x)=g(x)=(1+C)G(x)
 \end{gathered}
\end{equation}
for some fixed $\eta>0$.

Here is the full smoothing argument.  Put $y_j=1-x_j$,
$t=k^{-1}\sum_j y_j$, $h=F-G=t^k-\prod_jy_j$, and
$d=\max_{p,q}|y_p-y_q|$.  When $t>0$,
\begin{equation}\label{simp:smoothing-bounds}
 h\ge\frac{d^2t^{k-2}}{8k^2},\qquad
 |\partial_jh|\le k^{k-1}dt^{k-2},\qquad
 \frac{|\partial_jh\,\partial_lh|}{h}
                 \le8k^{2k}t^{k-2}\quad(h>0).
\end{equation}
To verify the first bound, set $w_j=y_j/t$ and
$V=\sum_j(w_j-1)^2$.  Since $0\le w_j\le k$ and
$\sum_jw_j=k$, the inequality
$\log w\le w-1-(w-1)^2/(2k^2)$ gives
$\prod_jy_j\le t^k e^{-V/(2k^2)}$.
Here $V\le k(k-1)$ and $V\ge d^2/(2t^2)$; apply
$1-e^{-u}\ge u/2$ for $0\le u\le1/2$.
The product-difference identity gives the second bound, since
$|y_j-t|\le d$ and all product factors are at most $kt$.
The third follows from the first two.  Zero factors follow by limits.

Choose $0<a<1$, put $b=\sqrt a$ and $\lambda=2/\log(1/a)$,
and let $\varphi(u)=\int_0^u\psi(v)\,dv$, where
\[
 \psi(v)=
 \begin{cases}
 1,&0\le v\le a,\\
 1-\lambda\log(v/a),&a<v<b,\\
 0,&v\ge b.
 \end{cases}
\]
Thus $0\le\varphi'\le1$, $0\le\varphi(u)\le\min(u,b)$,
and $\varphi''=-\lambda/u$ on $(a,b)$ and zero elsewhere.
Define
\[
 C=\frac{16\lambda k^{2k+1}}{k-1},\qquad
 f=F-\varphi(h)+CG,\qquad g=(1+C)G.
\]
Nonnegativity follows from $F-\varphi(h)\ge G$.
Also
\[
 \partial_jf=(1-\varphi'(h))\partial_jF
                   +(\varphi'(h)+C)\partial_jG,
\]
which gives the first derivative bounds.  Wherever second
derivatives exist,
\[
 \partial_{jl}f=(1-\varphi')\partial_{jl}F
    +\varphi'\partial_{jl}G
    -\varphi''\partial_jh\,\partial_lh+C\partial_{jl}G.
\]
The first two terms are nonpositive.  By
\eqref{simp:smoothing-bounds}, the third is at most
$8\lambda k^{2k}t^{k-2}$, whereas
$C\partial_{jl}G=-16\lambda k^{2k}t^{k-2}$.
Hence all second partials are nonpositive.  The functions are
$C^1$ and piecewise $C^2$, so this proves diminishing returns,
including across the joins.  Since $h=0$ on the diagonal,
uniform continuity gives $\eta>0$ with $d\le\eta$ implying
$h\le a$; there $\varphi(h)=h$.
Finally $f(e_j)\ge1-b$, and taking $a$ small ensures
$b,C<\zeta$.  This proves \eqref{simp:masked-properties}.
Concavity of $G$ in $\bar x$ also gives
\begin{equation}\label{simp:reference-sum}
 \sum_{i=1}^k g(x^i)\le k(1+C)\rho_k
 \quad\hbox{if}\quad \sum_i x^i_j\le1\quad(j=1,\ldots,k).
\end{equation}

\subsection{The construction, with rational values from the start}

Take $k$ divisible by three, $n=4k/3$, and $n$ reserved items $H$.
There are $k/3$ type-A agents, of value $|S\cap H|$, and $k$
type-B agents.  A further $ks$ items $M$ are uniformly secretly
equipartitioned into $X_1,\ldots,X_k$, each of size $s$.
Let $m=n+ks$, $\delta=s^{-1/2}$, and
$x_j(T)=|T\cap X_j|/s$.
For $v=f$ (the planted instance) or $v=g$ (the reference instance),
define the common type-B valuation using a fixed rational $0<\tau\le1$:
\begin{equation}\label{simp:rational-B}
 \begin{gathered}
 w_v(S)=\min\{1,\delta|S\cap H|+v(x(S\cap M))\},\qquad
 q_m(r)=\frac{2mr-r^2}{m^2},\quad \Delta=\frac{\tau}{10m^2},\\
 \widehat w_v(S)=\Delta\left\lfloor
           \frac{w_v(S)+\tau q_m(|S|)}{\Delta}\right\rfloor.
 \end{gathered}
\end{equation}
Continuous diminishing returns makes $v(x(T))$ submodular.
Adding the disjoint modular term preserves that property, as does
truncation: for a monotone submodular function $u$, the marginal of
$\min(1,u)$ is $\min\{u(S+e)-u(S),(1-u(S))^+\}$.
Thus $w_v$ is normalized, monotone, submodular, and at most one.

The floor operation also preserves these properties here, with
strict margins.  The minimum cardinality increment of $q_m$ is
$1/m^2$, and its elementary two-element submodular difference is
$2/m^2$.  Write $\widehat w_v=w_v+\tau q_m-e_S$, where
$0\le e_S<\Delta$.  Every rounded increment is at least
$\tau/m^2-\Delta>0$; every elementary submodular difference is at
least $2\tau/m^2-2\Delta>0$.  These elementary inequalities imply
all diminishing-returns inequalities by telescoping.  Moreover
\begin{equation}\label{simp:rational-error}
 w_v(S)-\Delta<\widehat w_v(S)
             \le w_v(S)+\tau\le1+\tau.
\end{equation}
Normalization is exact.  If $\tau=a'/b'$ in lowest terms, each
answer is $Ja'/(10b'm^2)$ with $Ja'\le10(a'+b')m^2$,
so its bit length is $O(\log m+\operatorname{bits}(\tau))$.
Type-A answers remain integers.

Every first singleton matching must select exactly $H$ for large $s$.
A type-A agent values its singletons in $H$ at one and all others
at zero.  For type B, the derivative bound in
\eqref{simp:masked-properties} gives residual singleton value at
most $(1+C)/s$ before rounding, while an $H$ singleton has value
$\delta$.  Their cardinality padding is identical, so the rounded
difference is greater than $\delta-(1+C)/s-\Delta>0$.
For instance $s\ge4(1+C)^2$ suffices.  A positive matching using
all of $H$ exists; any positive matching using a residual item
leaves an $H$ item unused and is strictly improved by exchanging
those two items.  A zero-valued matching is not optimal.

Give four reserved items to each type-A agent and one hidden block
to each type-B agent.  In the planted instance this yields
\begin{equation}\label{simp:planted-opt}
       \mathrm{OPT}\ge4^{1/4}(1-\zeta-\Delta)^{3/4}.
\end{equation}
In any reference allocation with exactly one reserved item per
agent, type-A values are one.  The disjoint type-B residual bundles
$T_i$ obey \eqref{simp:reference-sum}, so their average rounded
value is at most $\delta+\tau+(1+C)\rho_k$.
Arithmetic--geometric mean gives the uniform restricted bound
\begin{equation}\label{simp:reference-NSW}
       \mathrm{NSW}_{\rm ref}\le
                  [\delta+\tau+(1+C)\rho_k]^{3/4}.
\end{equation}

\subsection{Exact transcript masking and the quantifiers}

For each fixed $T\subseteq M$ and each hidden block, the variable
$Y_j=|T\cap X_j|$ samples $s$ items without replacement from $ks$,
with mean $|T|/k$.  Exposing these draws gives a conditional-mean
martingale whose increments have absolute value at most one: at
draw $r$ its two possible next conditional means differ by
$1-(s-r)/(ks-r)\in[0,1]$.
For a mean-zero increment $D$ with $|D|\le1$,
$\mathbb E e^{uD}\le\cosh u\le e^{u^2/2}$.
Multiplication and optimization of the exponential bound show
$\Pr\{|Y_j-|T|/k|>\eta s/2\}\le2e^{-\eta^2s/8}$.
A union bound therefore gives
\[
 \Pr\{\max_{p,q}|x_p(T)-x_q(T)|>\eta\}
                  \le2k e^{-\eta^2s/8}.
\]

Fix the algorithm's coins and run it on the reference oracle.
Every answer, including its cardinality padding and floor, depends
only on the two visible cardinalities $|S\cap H|,|S\cap M|$.
The entire reference query sequence and output are independent of
the hidden partition.  On a balanced query, \eqref{simp:masked-properties}
gives exactly equal inputs to the identical padding and floor in
\eqref{simp:rational-B}; the two oracle answers are exactly equal.
Union bound over the $q$ reference queries and the $k$ type-B
output bundles.  Except with probability
$2k(q+k)e^{-\eta^2s/8}$, the adaptive planted execution has the
identical transcript and output values.  Equations
\eqref{simp:planted-opt}--\eqref{simp:reference-NSW} then imply
\[
 \frac{\mathrm{OPT}}{\mathrm{NSW}_{\rm out}}\ge
 4^{1/4}\left(\frac{1-\zeta-\Delta}
                   {\delta+\tau+(1+C)\rho_k}\right)^{3/4}.
\]
Zero output welfare already implies failure.  Averaging over the
fixed coins proves the same assertion for randomized algorithms.

For a prescribed $\varepsilon$, first fix a sufficiently large $k$
divisible by three; then fix sufficiently small $\zeta$ and rational
$\tau$, and choose the masked pair with $C<\zeta$.  Its $\eta>0$
is now fixed.  Finally take $s\to\infty$.
Since $\rho_k\to1-e^{-1}$, the displayed ratio is greater than
$R_\star-\varepsilon$ for all large $s$.
Put $n_0=4k/3$ and $c=\eta^2/(32k)$.  For $q\le e^{cm}$,
the exceptional probability is at most
$\exp(-3\eta^2s/32+O_\varepsilon(1))=o(1)$.
This proves the theorem with $n_0$ fixed and short rational answers.

The conclusion also excludes expected welfare factor $1\le R<R_\star$.
Indeed, for $R<R'<R_\star$ and
$X=\mathrm{NSW}_{\rm out}/\mathrm{OPT}\in[0,1]$,
$\mathbb EX\ge1/R$ would contradict the vanishing success probability, since
\[
 \Pr\{X\ge1/R'\}\ge
               \frac{1/R-1/R'}{1-1/R'}>0.
\]
Together with the upper guarantee for every $R_\star+\varepsilon$, this
proves the exact infimum threshold under the stated framework restriction.
